% ============================================================================
\clearpage
\section*{v4.9.0-r24 Source-Only: Correction of the Three Discrepancies of r23 as Single-Frame Artefacts}
\addcontentsline{toc}{section}{v4.9.0-r24 Source-Only: Correction of the Three Discrepancies of r23 as Single-Frame Artefacts}
\label{sec:v490-r24-single-frame-correction}
% ============================================================================

\paragraph{Purpose.}
r23 recorded three discrepancies between the main body and Convention F: the sign of the matter offset, the position of the
tachyon, and the position of the photon.  On 2026-10-22 the author asked for a correction and pointed out that the main body
worked in one frame only, which may be where the discrepancies come from.  This layer shows that this is so for all three.  It
corrects them by typing, not by editing: the main body is not rewritten, and its statements receive the frame they were made
in.  The script is \texttt{llm/k63\_widersprueche\_ein\_frame.py}.  No prediction is made, and nothing is compared with data.

\noindent\textbf{Source-only status.}
This v4.9.0-r24 layer is source-only.  It does not build a PDF, does not assign a new DOI,
and does not authorize a public release.  The v4.9.0 Version DOI
\href{https://doi.org/10.5281/zenodo.23534954}{10.5281/zenodo.23534954} and the Concept DOI
\href{https://doi.org/10.5281/zenodo.19210728}{10.5281/zenodo.19210728} are unchanged by it.

\subsection*{v4.9.0-r24 the single frame of the main body}
\label{subsec:v490-r24-single-frame}

The main body describes everything from one frame and along one real axis.  It calls that frame the photon frame and measures
the offset \(\bsym\) from it.  Three features of that frame matter here.  First, it fixes an orientation on the axis.  Second, it
sees only real rapidities.  Third, it is the zero of its own scale line.  Each discrepancy of r23 comes from one of these.

\subsection*{v4.9.0-r24 the three corrections}
\label{subsec:v490-r24-corrections}

\paragraph{Theoremlet v4.9.0-r24.1 (sign).}
The sign of \(\beta\) in the sector definition of \S\ref{sec:unitary} (matter \(\beta<0\)) and in Theorem~\ref{thm:symmetric-frame}
(matter \(+\bsym\)) is an orientation code in the sense of v4.8.0-r01
(\S\ref{subsec:v480-r01-orientation-vs-signature}).  The two choices differ by the reflection \(z\mapsto-z\), which is charge
conjugation \(C\) (v4.9.0-r02).  \(C\) reverses \(\operatorname{Re}z\) and leaves \(s=z\cdot z\) unchanged.  The sector is fixed by the sign
of \(s\) (v4.8.0-r07), so both choices assign the same sectors.  The two passages use opposite orientation conventions.  That is
a conflict of notation, not of physics.

\paragraph{Correction.}
Read \(\beta\) of the sector definition and \(\pm\bsym\) of Theorem~\ref{thm:symmetric-frame} as orientation codes, and read
\(\bsym=|\beta|\ge0\) as the magnitude of the offset.  Matter and tachyon are told apart by the sign of \(s\).

\paragraph{Theoremlet v4.9.0-r24.2 (tachyon).}
The tachyon rapidity is \(-\eta+i\pi/2\), with \(\tanh(-\eta+i\pi/2)=-\coth\eta=-1/\bsym\).  Projected onto the single real axis of the
main body, dropping the imaginary part, it is \(-\eta\), that is \(\beta=-\bsym\).  This is the tachyon of
Theorem~\ref{thm:symmetric-frame}.  In its own frame the tachyon has \(\beta_T^2=1/\bsym^2\), which is the tachyon of Convention F.
Both describe the same tachyon, the first as seen along the axis of the reference frame and the second in its own frame.

\paragraph{Correction.}
Convention F receives a fourth entry, \(T|_R\), the tachyon seen from the reference frame of the main body, at \(-\bsym\).  The \(18\)
even statements of r23.1(a), among them \(\Omega_M\), \(g_{tt}=-\bsym^2\) and the Beltrami coefficient, hold for M and for \(T|_R\)
alike.  They are frame-free in the description of the main body.  In the tachyon's own frame they take the values at
\(1/\bsym^2\) (\(\Omega_M\to3.31\)).  The two values are different typings, and neither contradicts the other.

\paragraph{Theoremlet v4.9.0-r24.3 (photon).}
The frame that the main body calls the photon frame, \(\bsym=0\), is in the state disk the pure state \(|b|=1\), with entropy \(0\).
The main-body conversion factor to it is the ratio of Bloch lengths:
\[
\frac{|b_R|}{|b_M|}=\frac{1}{\operatorname{sech}2\eta}=\cosh2\eta=\gL .
\]
To the photon of Convention F, \(b=0\), the same factor would diverge.  The reference frame of the main body is therefore the
pure-state frame R, the zero of the scale line.  The photon P of Convention F is the maximally mixed state, reached for
\(\beta\to1\).

\paragraph{Proof sketch.}
r24.1: \(s(-z)=s(z)\) and \(\operatorname{Re}(-z)=-\operatorname{Re}z\), checked on \(1000\) random \(z\).  r24.2: \(\tanh(x+i\pi/2)=\coth x\).
r24.3: \(|b|=\operatorname{sech}2\eta\), equal to \(1\) at \(\eta=0\).  Numerically \(|b_R|/|b_M|=1.86444=\gL\).  \(\square\)

\paragraph{Correction.}
Read ``photon frame'' in the main body, where it denotes the zero of the offset or the target of a \(\gL^{-k}\) conversion, as the
reference frame R.  The nine S statements of r23 become R statements and hold there.  \(\Lambda=0\) and the vanishing frame term
hold in R, the frame where the doublet state is pure.  At the photon P the frame term \(\bsym^2\) and the ratio \(p_-/p_+\) are
maximal (\(1\)).  This also explains why a finite \(\gL\) points to \(|b|=1/\gL\) and not to \(b=0\) (r17): the conversion goes from M to
R, not from M to P.

\subsection*{v4.9.0-r24 Convention F completed}
\label{subsec:v490-r24-convention}

\begin{enumerate}[leftmargin=2em]
\item[(F5)] \emph{Reference frame.}  R is the pure state \(|b|=1\), \(\bsym=0\), \(\sigma=0\).  It is the frame of the main body.  The
conversion \(X^{(R)}=\gL^{-k}X^{(M)}\) of \S\ref{sec:phase-constants} goes from M to R.
\item[(F6)] \emph{Projection.}  \(T|_R\) is the tachyon seen from R along the real axis: \(\beta=-\bsym\), imaginary part dropped.  A
statement of the main body about the tachyon is a statement about \(T|_R\) unless it says otherwise.
\item[(F7)] \emph{Orientation.}  The sign of \(\beta\) is an orientation code.  The sector is fixed by the sign of \(s\).
\end{enumerate}
The r23 classification becomes: S \(\to\) R (valid in R); the \(18\) even statements are frame-free between M and \(T|_R\) and
matter-only in the tachyon's own frame; P keeps the three \(\beta\to1\) statements.

\subsection*{v4.9.0-r24 status in the three frames}
\label{subsec:v490-r24-frames}

\begin{center}
\small
\begin{tabular}{@{}>{\raggedright\arraybackslash}p{0.30\textwidth}>{\raggedright\arraybackslash}p{0.62\textwidth}@{}}
\toprule
result & M / T / P (and R, \(T|_R\))\\
\midrule
r24.1 sign & orientation code in every frame; sector by the sign of \(s\), unchanged under \(C\)\\
r24.2 tachyon & \(T|_R\) at \(-\bsym\), T at \(1/\bsym\); even statements equal in M and \(T|_R\), different in T\\
r24.3 photon & R is the pure state, P the maximally mixed state; S statements hold in R, not in P\\
\bottomrule
\end{tabular}
\end{center}

\subsection*{v4.9.0-r24 register entries}
\label{subsec:v490-r24-register}

\begin{center}
\small
\begin{tabular}{@{}>{\raggedright\arraybackslash}p{0.34\textwidth}>{\raggedright\arraybackslash}p{0.58\textwidth}@{}}
\toprule
statement & status after v4.9.0-r24\\
\midrule
sign of the matter offset & orientation code; the two conventions differ by \(C\); sector by \(s\) (r24.1)\\
tachyon of Thm.~\ref{thm:symmetric-frame} & \(T|_R\), the real projection of the tachyon onto the reference axis (r24.2)\\
photon frame of the main body & reference frame R, the pure state; conversions go from M to R (r24.3)\\
S statements of r23 & valid in R\\
\(18\) even statements of r23 & frame-free between M and \(T|_R\); in the tachyon's own frame at \(1/\bsym^2\)\\
discrepancies of r23 & resolved as single-frame artefacts; no edit of the main body\\
\bottomrule
\end{tabular}
\end{center}

\subsection*{v4.9.0-r24 ledger and guard update}
\label{subsec:v490-r24-ledger-guard}

\[
R_{\rm F}^{v490\text{-}r24}=(0,0,0)
\]
for (sign resolved, tachyon placement resolved, photon placement resolved).  The open component of r23, the choice between the
two tachyon placements, is closed: both are kept with their typings.

The layer spends no guard:
\[
R_{\rm glob}^{(v490\text{-}r24,+)}=(0,1,1,1,1),
\qquad
\Gamma_{\rm guard}^{v490\text{-}r24}=(1,1,1,1,1),
\qquad
\beta_4=1/256\ \text{unspent}.
\]
It does not imply exact finite \(C5\) zero closure, exact all-readout terminality,
\(C6/T2\), a global photon-clock/RH theorem, or global all-readout closure.

